최근 대규모 언어 모델(LLM)의 발전으로 인해 단일 모델의 추론 능력을 넘어, 특정 역할을 부여받은 여러 에이전트가 협력하여 복잡한 과업을 수행하는 멀티 에이전트 시스템(Multi-Agent System, MAS)에 대한 관심이 급증하고 있다. 특히 고도의 전문성과 구조적 완결성이 요구되는 문서 생성 작업의 경우, 설계, 집필, 검수, 수정으로 이어지는 파이프라인을 구축하여 각 단계의 전문성을 극대화하는 방식이 효과적이다. 이때 에이전트 간의 작업 흐름을 제어하고 상태를 관리하는 오케스트레이션(Orchestration) 메커니즘은 시스템의 전체적인 성능과 자원 효율성을 결정짓는 핵심 요소로 작용한다.

\subsection{오케스트레이션 제어 방식의 상충 관계}

멀티 에이전트 시스템의 제어 방식은 크게 결정론적 제어(Deterministic Control)와 확률론적 제어(Probabilistic Control)로 구분할 수 있다. 전자는 유향 비순환 그래프(DAG) 기반의 워크플로우와 같이 미리 정의된 경로를 따르는 방식이며, 후자는 AutoGPT와 같이 에이전트가 자율적으로 다음 행동을 결정하는 방식이다. 

고정된 그래프 구조를 따르는 Code 기반 오케스트레이션은 개발자가 정의한 엄격한 워크플로우에 따라 작동하므로 실행 경로가 예측 가능하다. 이론적으로 이러한 구조는 불필요한 추론 단계를 최소화하여 자원 효율성을 높일 수 있을 것으로 기대되며, 본 연구의 기초 실측 결과 Code 구조의 평균 실행 시간은 158.61초, 토큰 소비량은 21,440개로 나타났다. 그러나 이러한 결정론적 방식은 작업 도중 발생하는 예외 상황이나 LLM의 생성 결과에 따라 유연하게 경로를 수정할 수 없다는 경직성의 한계를 지닌다.

반면, 기존 연구들에 따르면 LLM이 직접 다음 단계를 결정하는 동적 라우팅(Dynamic Routing) 방식의 오케스트레이션은 작업의 맥락에 따라 실행 흐름(Execution Flow)을 실시간으로 조정할 수 있어 높은 유연성을 제공하는 것으로 알려져 있다. 하지만 이는 LLM의 확률적 특성으로 인해 제어 흐름의 예측 불가능성을 초래하며, 무한 루프(Loop) 발생 가능성 및 과도한 토큰 소비로 인한 비용 증가라는 잠재적 위험을 내포하고 있다. 이러한 경직성과 예측 불가능성의 대립은 시스템 설계자로 하여금 작업의 성격에 맞는 최적의 제어 구조를 선택해야 하는 과제를 안겨준다.

\subsection{연구의 목적 및 범위}

기존의 멀티 에이전트 연구들은 주로 개별 에이전트의 프롬프트 최적화나 추론 능력 향상에 집중해 왔으며, 오케스트레이션 구조 자체가 최종 생성물의 품질과 자원 소모에 미치는 영향을 정량적으로 비교한 연구는 상대적으로 부족하다. 따라서 본 연구의 목적은 서로 다른 세 가지 오케스트레이션 구조인 Code(고정 그래프), LLM(동적 라우팅), 그리고 Hybrid(계층적 제어) 방식이 문서 생성 성능에 미치는 영향을 정량적으로 비교 분석하여, 구조적 설계가 시스템 성능에 주는 실질적인 기여도를 규명하는 것이다. 특히, 산업 현장에서 요구되는 정형 보고서(예: 금형-사출 분석 보고서)와 같이 엄격한 형식이 중요한 작업과, 창의적 글쓰기와 같이 유연한 전개가 필요한 작업 사이에서 오케스트레이션 구조가 생성 품질과 자원 효율성 사이에서 어떠한 Trade-off 관계를 형성하는지 입증하고자 한다.

본 연구의 분석 범위는 다음과 같이 정의한다. 첫째, 독립변인으로 세 가지 오케스트레이션 구조를 설정하고, 통제변인으로 에이전트의 역할 정의, 사용 LLM 및 프롬프트 로직을 동일하게 유지하여 구조적 차이만을 검증한다. 둘째, 종속변인으로서 자원 효율성을 측정하기 위한 토큰 소비량과 실행 시간, 그리고 생성 효과성을 측정하기 위한 LLM-as-Judge 기반의 품질 점수와 재시도 횟수($\text{retry\_count}$)를 수집한다. 이를 통해 작업의 성격에 따라 어떤 오케스트레이션 구조가 최적인지를 판단할 수 있는 정량적 기준을 제시함으로써, 향후 효율적인 멀티 에이전트 시스템 설계의 가이드라인을 제공하는 데 기여하고자 한다.